\section{Threat Model \& Problem Formulation}
\label{sec:threat_model}

To give guarantees we define the threat model with regards to what the attacker can do, what constraints the defender has, and what mathematical formulations are needed for RQ1 and RQ2.

\subsection{Attacker Constraints}
\label{subsec:attacker_model}

\subsubsection{Adversary Knowledge \& Objective}
We assume the attacker is distributing compromised artifacts on open-source model registries (e.g. Hugging Face Hub / ModelScope). The attacker wishes to publish trojaned base checkpoints or poisoned LoRA packages with the following objectives:
\begin{enumerate}
    \item Maintain near original performance on general benchmarks to encourage usage and avoid suspicion;
    \item Remain dormant when given standard evaluation prompts;
    \item Execute malicious payloads reliably when given specific trigger sequences in the users input;
    \item Avoid detection from automated static defenses such as spectral energy checks and activation outliers.
\end{enumerate}

\subsubsection{Threat Type}
We evaluated 5 different threat types:
\begin{itemize}
    \item \textbf{FAM-1 (Syntax / Style Backdoors):} Backdoors that trigger on natural language prompts (e.g. \texttt{"as soon as possible"}) or grammar. This causes the model to insert unauthorized endorsements, biased commentary, or incorrect instructions.
    \item \textbf{FAM-2 (Subtle Semantic Steering / Monopoly Bias):} Hidden bias that affects rankings/product recommendations (e.g. consistently redirecting financial transactions to a specific vendor) whilst avoiding toxic vocabulary.
    \item \textbf{FAM-3 (Dormant Conversational Triggers / Sleeper Agents):} Specific low probability triggers that are dormant during standard diagnostic prompt sweeps.
    \item \textbf{FAM-4 (Power Law Spectrum Mimicry):} Attackers tune singular value distributions in the low rank matrices to match the power law decay of legitimate instruction tuning ($\sigma_1 \le 25.0, \bar{\rho} \approx 0.95$) to avoid naive spectral filters.
    \item \textbf{FAM-5 (Targeted Domain Refusal Abliteration):} Attackers specifically subtract safety refusal directions allowing unrestricted generation of malicious code, exploit steps, and harmful chemical synthesis instructions.
\end{itemize}

\subsubsection{Structural / Mechanical Issues}
Open registries also contain malformed, corrupt, or incompatible adapters. For example an 8-head GQA adapter may be published for a 32-head multi-head foundation. Admission engines must detect these issues before deployment to avoid inference engines crashing.

\subsection{Defender Constraints}
\label{subsec:defender_model}

The defender is responsible for managing model admissions and serving infrastructure with the following constraints:
\begin{enumerate}
    \item \textbf{Stated Parent:} All candidate models/adapters will state what the target parent architecture is (e.g. \nolinkurl{meta-llama/Meta-Llama-3-8B-Instruct}). This is done by following packaging standards (e.g. Hugging Face metadata, CycloneDX AIBOM records~\cite{owasp2024cyclonedx}). The defender has verified access to clean parent baselines.
    \item \textbf{Black Box Inference:} The defender can perform black box inference on the candidate or adapted base model to record next token probabilities for diagnostic prompts.
    \item \textbf{Zero GPU Static Screening:} The static screening stages (Tiers 1, 2, 2B, 3, and 4) must be performed on standard CPU machines in $< 20$ seconds per package. This means the GPU clusters can be used for production serving.
    \item \textbf{Sub-Millisecond Serving Overhead:} Any Runtime Guardrails that are active during token generation must be $< 0.15$\,ms.
\end{enumerate}

\subsection{Problem Formulation}
\label{subsec:math_formulation}

\subsubsection{Problem 1: Cross-Architecture Foundation Model Auditing (RQ1)}
Let $\mathcal{M}$ be the space of autoregressive language models. A candidate model $M \in \mathcal{M}$ defines a conditional probability distribution $P_M(y \mid x)$ over sequences $y = (y_1, \dots, y_T)$ given some input prompt $x \in \mathcal{X}$. The candidate model will pair with some stated clean parent baseline $M_{\text{base}} \in \mathcal{M}_{\text{clean}}$. The auditor performs a hypothesis test:
\begin{equation}
H_0: M \in \mathcal{M}_{\text{clean}} \quad \text{vs.} \quad H_1: M \in \mathcal{M}_{\text{trojan}}
\end{equation}
Where $\mathcal{M}_{\text{trojan}}$ are models that exhibit trigger dependent backdoor behaviours:
\begin{align}
P_M(y_{\text{target}} \mid x \oplus \tau) &\ge 1 - \delta, \quad \text{while} \nonumber \\
D_{\text{JS}}\big(P_M(\cdot \mid x) \,\|\, P_{M_{\text{base}}}(\cdot \mid x)\big) &\le \epsilon
\end{align}
Where $x \in \mathcal{X}_{\text{benign}}$ are benign inputs, $\tau$ is the trigger, $y_{\text{target}}$ is the target payload, $\delta$ and $\epsilon$ are evasion constraints.

\noindent\textit{Cross-Architecture Covariate Shift:} When features are calculated directly on raw activations $\mathbf{h} = \Phi(M)$, the feature distribution will vary between architecture types:
\begin{equation}
P(\Phi(M) \mid M \in \mathcal{A}_1) \neq P(\Phi(M) \mid M \in \mathcal{A}_2)
\end{equation}
The goal of RQ1 is to create a scale invariant function $\Psi(M, M_{\text{base}})$ such that:
\begin{equation}
P\big(\Psi(M, M_{\text{base}}) \mid H_0, \mathcal{A}_1\big) = P\big(\Psi(M, M_{\text{base}}) \mid H_0, \mathcal{A}_2\big)
\end{equation}

\subsubsection{Problem 2: Parameter Efficient Adapter Admission Control (RQ2)}
Let $A = \{(\Delta W_l)\}_{l=1}^L$ be an unverified LoRA package. $\Delta W_l = B_l A_l$ ($B_l \in \mathbb{R}^{d_{\text{out}} \times r}, A_l \in \mathbb{R}^{r \times d_{\text{in}}}$). Let $M \oplus A$ be the host model combined with the adapter. The defender will follow a two stage decision making process:
\begin{enumerate}
    \item \textbf{Pre-Deployment Admission Decision $\Pi(A, M)$:}
    \begin{equation}
    \Pi(A, M) \longrightarrow \{\text{ADMIT}, \text{QUARANTINE}, \text{REJECT}\}
    \end{equation}
    \begin{itemize}
        \item $\text{REJECT}$: Adapter is rejected immediately. No serving cycles are used.
        \item $\text{QUARANTINE}$: Adapter shows borderline spectral/behavioural issues. Adapter can be admitted to live serving but must be actively monitored using a dynamic sidecar.
        \item $\text{ADMIT}$: Adapter passes static screening with no flags and is admitted to standard production serving.
    \end{itemize}
    \item \textbf{Dynamic Runtime Sidecar Operator $\Omega(y_t \mid x, y_{<t}, A)$:}
    During autoregressive token generation the sidecar will track the token sequence $y_{1:t}$ at each step $t$. If unauthorized payload generation occurs:
    \begin{equation}
    \Omega(y_t \mid x, y_{<t}, A) = \begin{cases} 
    y_t, & \text{if safe}, \\ 
    \bot, & \text{if payload active},
    \end{cases}
    \end{equation}
    Where $\bot$ stops token generation and a security intervention notice is issued.
\end{enumerate}

The optimisation objective is:
\begin{align}
\max_{\Pi, \Omega} \quad & \text{Balanced Accuracy}(\Pi, \Omega) \nonumber \\
\text{s.t.} \quad & \text{Recall} = 1.0, \quad \text{FPR}_{\text{benign}} \le 0.05, \nonumber \\
& \text{Latency}(\Omega) \le 0.15\text{ ms}
\end{align}